\subsubsection{Receptive fields under unsupervised learning}
Receptive field (RF) connectivity outperforms random connectivity on every dataset, but
none of the unsupervised plasticity rules we tested improves on the RFs it starts from.
Averaged over the four datasets, \textit{Frozen} reaches $82.9\%$ accuracy, against
$80.5\%$ for \textit{Base} and $77.9\%$ for \textit{Random}
(Figure~\ref{fig:phase_1_predicted}). The beta-binomial GLMM
(Table~\ref{tab:glmm_phase1}, Appendix~\ref{sec:app:glmm_unsup}) takes \textit{Base} on
MNIST as its reference, so each configuration coefficient is the logit-scale effect of one
change on MNIST, and each interaction is how that effect shifts on another dataset.

\begin{figure}[h]
    \centering
    \includegraphics[width=0.7\linewidth]{fig_phase1_average_BW.pdf}
    \caption{Predicted linear-probe accuracy of the unsupervised model, averaged over the
    four datasets, for the six conditions that each make one change to \textit{Base}
    (dashed line). Boxes give the uncertainty of the model's mean prediction, simulated
    from the fitted fixed effects of the beta-binomial GLMM, and open circles the observed
    accuracy of each of the five seeds, averaged over the same four datasets.
    \textit{Base} is the oriented receptive fields condition under trace-STDP with
    recurrent excitation and fixed lateral inhibition; each other condition adds ($+$),
    removes ($-$), or replaces (unsigned) one component (Section~\ref{sec:exp1}).}
    \label{fig:phase_1_predicted}
\end{figure}

Disabling plasticity has the largest effect of any single change on MNIST
($\hat{\beta}_3 = 0.290$, $p<.001$), and \textit{Frozen} stays ahead of \textit{Base} on
KMNIST and notMNIST; only on FMNIST does a negative interaction
($\hat{\beta}_{12} = -0.267$, $p<.001$) cancel the advantage. Neither alternative learning
rule differs from trace-STDP on MNIST, with $\hat{\beta}_5 = 0.041$ ($p = .224$) for
\textit{Triplet} and $\hat{\beta}_6 = 0.003$ ($p = .908$) for \textit{$+$iSTDP}. \textit{Triplet} gains
on KMNIST and notMNIST through positive interactions ($\hat{\beta}_{20} = 0.105$ and
$\hat{\beta}_{26} = 0.173$, $p \leq .002$) but loses on FMNIST, whereas none of the
\textit{$+$iSTDP} interactions differs from zero.

Removing lateral inhibition is the more harmful connectivity change, with the largest
penalty of any condition on MNIST ($\hat{\beta}_4 = -0.203$, $p<.001$), whereas removing
recurrent excitation slightly improves accuracy on MNIST ($\hat{\beta}_2 = 0.082$,
$p = .004$), which indicates that recurrence contributes little to the representation.
Replacing the RFs with random connectivity costs comparatively little on MNIST
($\hat{\beta}_1 = -0.072$, $p = .013$), but all of its interactions are negative, so the
penalty grows on the other datasets and \textit{Random} ends up the least accurate
configuration. The seed intercept is singular and the dispersion very large
($\hat{\phi} = 1.8\times10^{5}$), so the variation between seeds stays within binomial
sampling noise and these effects are reproducible across seeds.

Taken together, the weight structure supplied by the receptive fields is worth more than
anything unsupervised plasticity adds to it. Receptive field connectivity outperforms random
wiring, and every rule we tested erodes that advantage, \textit{Triplet} least of all. These
findings shape the next stage of our experiments, in which we keep lateral inhibition,
remove recurrent excitation, and redesign the network architecture around the supervised
learning rule R-Hebb.
